\section{Related Work\label{sec:related}}

\noindent\textbf{Deeper architectures.} %
%
VGGNets~\cite{simonyan_iclr2015vgg} and Inception models~\cite{szegedy_cvpr2015googlenet} showed that increasing the depth of a network could significantly increase the quality of representations that it was capable of learning. By regulating the distribution of the inputs to each layer, Batch Normalization (BN)~\cite{ioffe_icml2015bn} added stability to the learning process in deep networks and produced smoother optimisation surfaces~\cite{santurkar_nips2018}. Building on these works, ResNets demonstrated that it was possible to learn considerably deeper and stronger networks through the use of identity-based skip connections~\cite{he_cvpr2016resnet, he_eccv2016preact}. Highway networks \cite{srivastava_nips2015highway} introduced a gating mechanism to regulate the flow of information along shortcut connections. Following these works, there have been further reformulations of the connections between network layers \cite{chen_nips2017dpn, huang_cvpr2017dns}, which show promising improvements to the learning and representational properties of deep networks.


An alternative, but closely related line of research has focused on methods to improve the functional form of the computational elements contained within a network.  Grouped convolutions have proven to be a popular approach for increasing the cardinality of learned transformations \cite{ioannou_cvpr2017roots, xie_cvpr2017resnext}. More flexible compositions of operators can be achieved with multi-branch convolutions \cite{ioffe_icml2015bn, szegedy_cvpr2015googlenet, szegedy_cvpr2016inceptionv3, szegedy_iclrw2016inceptionv4}, which can be viewed as a natural extension of the grouping operator. In prior work, cross-channel correlations are typically mapped as new combinations of features, either independently of spatial structure \cite{max_bmvc2014lowrank, chollet_cvpr2017xception} or jointly by using standard convolutional filters \cite{lin_iclr2014nin} with $1\times1$ convolutions. Much of this research has concentrated on the objective of reducing model and computational complexity, reflecting an assumption that channel relationships can be formulated as a composition of instance-agnostic functions with local receptive fields.  In contrast, we claim that providing the unit with a mechanism to explicitly model dynamic, non-linear dependencies between channels using global information can ease the learning process, and significantly enhance the representational power of the network.

\vspace{5pt}
\noindent\textbf{Algorithmic Architecture Search.} Alongside the works described above, there is also a rich history of research that aims to forgo manual architecture design and instead seeks to learn the structure of the network automatically.  Much of the early work in this domain was conducted in the neuro-evolution community, which established methods for searching across network topologies with evolutionary methods \cite{miller_1989,stanley_2002}. While often computationally demanding, evolutionary search has had notable successes which include finding good memory cells for sequence models \cite{bayer_icann2009,jozefowicz_icml2015} and learning sophisticated architectures for large-scale image classification \cite{xie_iccv2017,real_icml2017,real_2018a}.  With the goal of reducing the computational burden of these methods, efficient alternatives to this approach have been proposed based on Lamarckian inheritance~\cite{elsken_2018mul} and differentiable architecture search~\cite{liu_2018darts}.

By formulating architecture search as hyperparameter optimisation, random search \cite{bergstra_jmlr2012} and other more sophisticated model-based optimisation techniques \cite{liu_eccv2018progressive,negrinho_2017deeparchitect} can also be used to tackle the problem.  Topology selection as a path through a fabric of possible designs \cite{saxena_nips2016} and direct architecture prediction \cite{brock_iclr2018,baker_iclrw2018} have been proposed as additional viable architecture search tools. Particularly strong results have been achieved with techniques from reinforcement learning \cite{baker_iclr2017,zoph_iclr2017,zoph_cvpr2018,liu_iclr2018deepmind,pham_icml2018}.  SE blocks can be used as atomic building blocks for these search algorithms, and were demonstrated to be highly effective in this capacity in concurrent work \cite{tan_2018mnasnet}.


\vspace{5pt}
\noindent\textbf{Attention and gating mechanisms.}  Attention can be interpreted as a means of biasing the allocation of available computational resources towards the most informative components of a signal \cite{olshausen_1993, itti_1998model, itti_2001, larochelle_nips2010, mnih_nips2014, vaswani_nips2017}.  Attention mechanisms have demonstrated their utility across many tasks including sequence learning \cite{bluche_nips2016text, miech_2017contextgate}, localisation and understanding in images \cite{cao_iccv2015twice, max_nips2015stn}, image captioning \cite{xu_icml2015, chen_cvpr2017sca} and lip reading \cite{chung_cvpr2017lip}. In these applications, it can  be incorporated as an operator following one or more layers representing higher-level abstractions for adaptation between modalities.   Some works provide interesting studies into the combined use of spatial and channel attention \cite{wangfei_cvpr2017resattent, woo_eccv2018cbam}. Wang et al. \cite{wangfei_cvpr2017resattent} introduced a powerful trunk-and-mask attention mechanism based on hourglass modules \cite{newell_eccv2016hourglass} that is inserted between the intermediate stages of deep residual networks. By contrast, our proposed SE block comprises a lightweight gating mechanism which focuses on enhancing the representational power of the network by modelling channel-wise relationships in a computationally efficient manner.
